\scanpage{71}
\mspage{31}
satisfaisant les conditions a) c), \inserted{d)} de p. 28.
A-t-on aussi la condition b) de p. 28
\[
\Sigma_{ik}\cap\Sigma_{jk}=\varnothing\quad\hbox{si }i,j<k,\quad
i,j\ \hbox{non comparables ?}
\]
Si on avait $x\in\Sigma_{ik}\cap\Sigma_{jk}\subset\Sigma_k$,
alors on aurait deux éléments $\unsure{x_i}$ et $x_j$ de $\Sigma$,
au dessus de $i$ et $j$ respectivement, qui sont $<x$
(au dessus de $k$), donc (comme $\Sigma_{<x}$ est tot\textsuperscript{t} ordonné)
\deleted{on aurait} $x$ et $y$ seraient comparables, donc aussi
$i$ et $j$, absurde.

\uline{Conséquence} les structures ``déploiements
équisinguliers de stratifications, indexés par $I$'',
décrites \deleted{dans} la p. 28, équivalent à la structure suivante :

\margin{\unsure{et|si} les comp. connexes des $\Sigma_i$ sont de dim finie.}

\textcircled{a} Un \deleted{$\Sigma$} espace top. ordonné $\Sigma$,
satisfaisant aux conditions a) b) c) et 1°) de la page 27,
et \uline{de plus} : la condition que, \deleted{pour tout}
posant $\widetilde I=\pi_0(\Sigma)$, on ait, pour
$i,j\in\widetilde I$, $i<j$, que le morphisme d'immersion
propre bordant $\Sigma_{ij}\xrightarrow{b_{ij}}\Sigma_j$ soit
\uline{mono}

\textcircled{b} Une application strictement croissante
(\unsure{pas néc.} surjection)
\[
\widetilde I\longrightarrow I
\]
\mspage{70}
